\unit{Bài trình bày 149 — Các định lý đối ngẫu cho các bó đại số coherent}
\sid{149}
\src{169}
\seminar{(Tháng 5 năm 1957)}
\paperhead{CÁC ĐỊNH LÝ ĐỐI NGẪU CHO CÁC BÓ ĐẠI SỐ COHERENT}%
  {do Alexandre GROTHENDIECK}

Những kết quả sau đây, được gợi ý bởi \fquote{định lý đối ngẫu đại số}
của Serre, đã được tìm ra vào mùa đông năm 1955 và mùa đông năm 1956. Chúng được chứng minh
rất đơn giản nhờ những kết quả khá sơ cấp về đối đồng điều của các không gian
xạ ảnh [3], cùng với việc sử dụng rộng rãi đại số đồng điều của
Cartan--Eilenberg theo dạng trình bày trong [2].

\fsection{1}{Các Ext của các bó môđun ({[2]}, Chương 3 và 4)}

Cho $X$ là một không gian tôpô được trang bị một bó $\Osh$ các vành có đơn vị
(không nhất thiết giao hoán). Ta xét phạm trù Abel
$\Ccat^{\Osh}$ của các bó $\Osh$-môđun, cũng được gọi là $\Osh$-Môđun.
Mọi vật của phạm trù này đều có một phân giải nội xạ; nhờ đó
có thể định nghĩa trong phạm trù này các hàm tử $\operatorname{Ext}$ với những tính chất
hình thức quen thuộc. Cụ thể hơn, để tránh nhầm lẫn, ta
ký hiệu $\operatorname{Hom}_{\Osh}(X;A,B)$, hoặc đơn giản là
$\operatorname{Hom}(X;A,B)$, là nhóm Abel các $\Osh$-đồng cấu từ $A$
vào $B$, trong khi $\Homsh_{\Osh}(A,B)$ ký hiệu bó các mầm
đồng cấu từ $A$ vào $B$ ($A,B\in\Ccat^{\Osh}$). Với
$A\in\Ccat^{\Osh}$ cố định, ta định nghĩa các hàm tử $h_A$ và $\hsh_A$ nhận giá trị tương ứng
trong phạm trù $\Ccat$ các nhóm Abel và phạm trù
$\Ccat^{\Zsh}$ các bó nhóm Abel trên $X$, bằng các công thức:
\begin{equation}\tag{1.1}
  h_A(B)=\operatorname{Hom}_{\Osh}(X;A,B),
  \qquad
  \hsh_A(B)=\Homsh_{\Osh}(A,B).
\end{equation}

Các hàm tử $h_A$ và $\hsh_A$ là những hàm tử hiệp biến khớp trái;
ta xét các hàm tử dẫn xuất phải của chúng, được ký hiệu tương ứng bởi
$\operatorname{Ext}_{\Osh}^{p}(X;A,B)$ và $\Extsh_{\Osh}^{p}(A,B)$. Do đó,
theo định nghĩa,
\begin{equation}\tag{1.2}
\left\{
\begin{aligned}
\operatorname{Ext}_{\Osh}^{p}(X;A,B)
  &= (R^{p}h_A)(B)
   = H^{p}\bigl(\operatorname{Hom}_{\Osh}(X;A,C(B))\bigr),\\
\Extsh_{\Osh}^{p}(A,B)
  &= (R^{p}\hsh_A)(B)
   = H^{p}\bigl(\Homsh_{\Osh}(A,C(B))\bigr),
\end{aligned}
\right.
\end{equation}

trong đó ký hiệu $R^{p}$ chỉ việc chuyển sang các hàm tử dẫn xuất phải, còn
$C(B)$ ký hiệu một phân giải nội xạ tùy ý của $B$ trong $\Ccat^{\Osh}$.
Ký hiệu
\[
  \Gamma:\Ccat^{\Zsh}\longrightarrow\Ccat
\]
là hàm tử \fquote{tiết diện}. Nhắc lại rằng các hàm tử dẫn xuất phải của nó được
ký hiệu bởi
\[
  B\longmapsto H^{p}(X,B).
\]


\src{170}
\begin{equation}\tag{1.3}
  H^{p}(X,B)=(R^{p}\Gamma)(B)=H^{p}\bigl(\Gamma(C(B))\bigr).
\end{equation}
Hiển nhiên, ta có $h_A=\Gamma\hsh_A$; mặt khác, ta kiểm tra được rằng $\hsh_A$
biến các vật nội xạ thành các vật $\Gamma$-acyclic. Từ đó, theo cách quen thuộc,
ta suy ra:

\result{MỆNH ĐỀ}{1} Với mọi $\Osh$-môđun $A$, tồn tại một hàm tử
phổ đối đồng điều trên $\Ccat^{\Osh}$, hội tụ tới hàm tử phân bậc
$\bigl(\operatorname{Ext}_{\Osh}^{*}(X;A,B)\bigr)$ và có số hạng ban đầu là
\begin{equation}\tag{1.4}
  E_{2}^{p,q}(A,B)=H^{p}\bigl(X,\Extsh_{\Osh}^{q}(A,B)\bigr).
\end{equation}
Từ đó, ta thu được các \fquote{đồng cấu biên} và một dãy khớp năm số hạng
mà ta sẽ không viết ra.

\result{HỆ QUẢ}{1} Nếu $A$ đẳng cấu địa phương với $\Osh^{n}$, thì ta có
các đẳng cấu chính tắc
\begin{equation}\tag{1.5}
  \operatorname{Ext}_{\Osh}^{p}(X;A,B)
  \overset{\sim}{\longleftarrow}
  H^{p}\bigl(X,\Homsh_{\Osh}(A,B)\bigr)
\end{equation}
(do các đồng cấu biên của dãy phổ cho). Đặc biệt, ta có
một đẳng cấu chính tắc
\begin{equation}\tag{1.6}
  \operatorname{Ext}_{\Osh}^{p}(X;\Osh,B)=H^{p}(X,B).
\end{equation}

Để sử dụng các kết quả này, cần mô tả tường minh các
$\Extsh_{\Osh}^{p}(A,B)$. Các hàm tử này được tính địa phương;
nghĩa là, nếu $U$ là một tập mở của $X$, thì
\[
  \Extsh_{\Osh}^{p}(A,B)|U
  =\Extsh_{\Osh|U}^{p}(A|U,B|U),
\]
vì phép hạn chế lên $U$ của một $\Osh$-môđun nội xạ
là một $(\Osh|U)$-môđun nội xạ. Hơn nữa, với $x\in X$ cố định, ta có các đồng cấu có tính hàm tử
\begin{equation}\tag{1.7}
  \Homsh_{\Osh}(A,B)_x
  \longrightarrow
  \operatorname{Hom}_{\Osh_x}(A_x,B_x),
\end{equation}
chúng mở rộng duy nhất thành một đồng cấu giữa các hàm tử đối đồng điều
theo $B$:
\begin{equation}\tag{1.8}
  \Extsh_{\Osh}^{p}(A,B)_x
  \longrightarrow
  \operatorname{Ext}_{\Osh_x}^{p}(A_x,B_x).
\end{equation}

\result{MỆNH ĐỀ}{2} Nếu $A$, trong một lân cận của $x$, đẳng cấu với
đối hạch của một đồng cấu $\Osh^{m}\longrightarrow\Osh^{n}$, thì (1.7) là một
đẳng cấu với mọi $p$.
Điều này đúng,
nói riêng, nếu $A$ là một $\Osh$-môđun coherent [3].

\src{171}
\result{MỆNH ĐỀ}{3} Cho $L_{*}=(L_i)$ là một phân giải trái của
$\Osh$-môđun $A$ bởi các $\Osh$-môđun, mỗi môđun đẳng cấu địa phương với
một $\Osh^{n}$. Khi đó $\Extsh_{\Osh}^{*}(A,B)$ được đồng nhất với
$H^{*}\bigl(\Homsh_{\Osh}(L_{*},B)\bigr)$, và
$\operatorname{Ext}_{\Osh}^{*}(X;A,B)$ được đồng nhất với siêu đối đồng điều của $X$
đối với phức $\Homsh_{\Osh}(L_{*},B)$.

Chứng minh là chuẩn: ta xét song phức
$\Homsh_{\Osh}(L_{*},C(B))$, trong đó $C(B)$ là một phân giải nội xạ của $B$, cùng các
đồng cấu tự nhiên từ $\Homsh_{\Osh}(L_{*},B)$ và từ
$\Homsh_{\Osh}(A,C(B))$ vào song phức này.

Cuối cùng, lưu ý rằng hai hàm tử $\operatorname{Ext}$ được đưa vào
trong (1.2) không chỉ là các hàm tử đối đồng điều theo $B$, mà còn là các
song hàm tử đối đồng điều, hiệp biến theo $B$ và phản biến theo $A$.

\fsection{2}{Luật hợp thành trong các Ext}

Các kết quả của mục này do CARTIER và YONEDA tìm ra độc lập; xem
một bài trình bày của CARTIER [1] để biết chi tiết. Cho $\Ccat$ là một phạm trù Abel.
Cho $K$ và $L$ là hai vật phân bậc trong $\Ccat$. Ta ký hiệu
$\operatorname{Hom}(K,L)$ là nhóm Abel phân bậc mà thành phần bậc
$n$ gồm các đồng cấu thuần nhất bậc $n$ từ $K$ vào $L$ (nghĩa là
các hệ $(u_i)$ đồng cấu $K^{i}\longrightarrow L^{i+n}$).

Nếu $K$ và $L$ là các phức (với các toán tử vi phân bậc $+1$, để
cụ thể hóa), ta trang bị cho $\operatorname{Hom}(K,L)$ toán tử
vi phân cho bởi
\begin{equation}\tag{2.1}
  \delta(u)=du+(-1)^{n+1}ud,\qquad\text{trong đó }n=\deg(u),
\end{equation}
qua đó biến nó thành một phức với toán tử vi phân bậc $+1$. Các chu trình
bậc $n$ là những toán tử bậc $n$ phản giao hoán với $u$ (xét như
các toán tử thuần nhất). Khi đó ta có thể xét
$H^{*}\bigl(\operatorname{Hom}(K,L)\bigr)$; đây là một bất biến của các kiểu
đồng luân của $K$ và $L$, và ta có thể ký hiệu nó là $H^{*}(K,L)$.

Nếu có thêm một phức thứ ba $M$, thì phép hợp thành các đồng cấu
xác định một phép ghép cặp
\[
  \operatorname{Hom}(K,L)\times\operatorname{Hom}(L,M)
  \longrightarrow\operatorname{Hom}(K,M)
\]
tương thích với các toán tử vi phân; do đó, khi chuyển sang đối đồng điều,
ta được các phép ghép cặp
\begin{equation}\tag{2.2}
  H^{*}(K,L)\times H^{*}(L,M)\longrightarrow H^{*}(K,M),
\end{equation}
được ký hiệu $(u,v)\longmapsto vu$. Các phép ghép cặp này thỏa mãn tính
kết hợp hiển nhiên. Đặc biệt, $H^{*}(K,K)$ là một vành phân bậc kết hợp
có đơn vị, và $H^{*}(K,L)$ (tương ứng $H^{*}(L,K)$) là một môđun phân bậc phải
(tương ứng trái) trên vành này, v.v. Ở bậc $0$, (2.2) thu về
phép hợp thành các đồng cấu của phức. Cuối cùng, một dãy khớp các
phức
\[
  0\longrightarrow K'\longrightarrow K\longrightarrow K''\longrightarrow 0,
\]

\src{172}
sao cho $K'^{i}$ được đồng nhất với một hạng tử trực tiếp của $K^{i}$ với mọi $i$,
dẫn đến một dãy khớp các phức nhóm $\operatorname{Hom}(K'',L)$, v.v., và do đó dẫn đến một
toán tử nối
\[
  H^{i}(K',L)\longrightarrow H^{i+1}(K'',L).
\]
Các toán tử nối tương ứng với một dãy khớp trong $L$ cũng được định nghĩa tương tự. Các
phép ghép cặp (2.2) tương thích, theo nghĩa thông thường, với các toán tử nối này.

Giả sử bây giờ $\Ccat$ là một phạm trù trong đó mọi vật $A$ của $\Ccat$
đều có một phân giải nội xạ $C(A)$. Khi đó, dùng một trong nhiều biến thể của định lý song phức, ta có
\[
  H^{*}\bigl(C(A),C(B)\bigr)
  =H^{*}\bigl(\operatorname{Hom}(C(A),C(B))\bigr)
\]
là đẳng cấu chính tắc với
\[
  H^{*}\bigl(\operatorname{Hom}(A,C(B))\bigr)
  =\operatorname{Ext}^{*}(A,B).
\]
Các toán tử nối được xét ở trên cho các đồng cấu nối của
$\operatorname{Ext}$. Hơn nữa, các phép ghép cặp (2.2) ở đây cho
các phép ghép cặp kết hợp, tương thích với các đồng cấu nối:
\begin{equation}\tag{2.3}
  \operatorname{Ext}^{*}(A,B)\times\operatorname{Ext}^{*}(B,C)
  \longrightarrow\operatorname{Ext}^{*}(A,C).
\end{equation}
Đặc biệt, $\operatorname{Ext}^{*}(A,A)$ là một vành phân bậc kết hợp
có đơn vị, v.v. (Tương tự, người ta chứng minh rằng các hàm tử
$\operatorname{Ext}$ tác động trên các hàm tử dẫn xuất của một hàm tử
bất kỳ; ở đây chúng ta sẽ không dùng đến sự kiện này.)

Trong trường hợp phạm trù đang xét là phạm trù $\Ccat^{\Osh}$ gồm các
$\Osh$-môđun trên $X$, do đó ta thu được các phép ghép cặp
\begin{equation}\tag{2.4}
  \operatorname{Ext}_{\Osh}^{p}(X;A,B)
  \times\operatorname{Ext}_{\Osh}^{q}(X;B,C)
  \longrightarrow\operatorname{Ext}_{\Osh}^{p+q}(X;A,C),
\end{equation}
có thể được viết tường minh như đã chỉ ra ở trên. Hơn nữa, bằng những triển khai tương tự,
trong đó thay phạm trù các nhóm Abel bằng phạm trù các bó nhóm Abel trên $X$, và thay các hàm tử
$\operatorname{Hom}$ bằng các hàm tử $\Homsh$, ta cũng định nghĩa được các phép ghép cặp có cùng
các tính chất hình thức và lần này mang \fquote{tính địa phương}:
\begin{equation}\tag{2.5}
  \Extsh_{\Osh}^{p}(A,B)\times\Extsh_{\Osh}^{q}(B,C)
  \longrightarrow\Extsh_{\Osh}^{p+q}(A,C).
\end{equation}
Các phép ghép cặp sau được làm rõ hơn nếu nhận thấy rằng các đồng cấu (1.8) tương thích
với các phép ghép cặp giữa các nhóm $\operatorname{Ext}$.

\src{173}
Cuối cùng, ta nhắc lại rằng còn có một cấu trúc nhân giữa các hàm tử
$H^{p}(X,A)$ (tích cup). Khi đó, ta thấy rằng các dãy phổ của
Mệnh đề 1 tương thích với các cấu trúc nhân; cụ thể hơn,
ta có một phép ghép cặp từ dãy phổ $E(A,B)$ và dãy phổ
$E(B,C)$ vào dãy phổ $E(A,C)$; tại đích hội tụ, phép này
thu về phép ghép cặp giữa các nhóm $\operatorname{Ext}$ toàn cục, còn tại
số hạng đầu thì đó là phép ghép cặp được suy ra, ở vế phải của (1.4), từ
tích cup và các phép ghép cặp của những nhóm $\operatorname{Ext}$ địa phương. Đặc biệt, từ đó suy ra
rằng các \fquote{đồng cấu biên}
\begin{equation}\tag{2.6}
  \operatorname{Ext}_{\Osh}^{n}(X;A,B)
  \longrightarrow H^{0}\bigl(X,\Extsh_{\Osh}^{n}(A,B)\bigr)
\end{equation}
và
\begin{equation}\tag{2.7}
  H^{n}\bigl(X,\Homsh_{\Osh}(A,B)\bigr)
  \longrightarrow\operatorname{Ext}_{\Osh}^{n}(X;A,B)
\end{equation}
tương thích với các cấu trúc nhân. Vì vậy, nếu chỉ xét
các bó đẳng cấu địa phương với một $\Osh^{m}$, thì điều này mô tả hoàn toàn
phép hợp thành giữa các nhóm $\operatorname{Ext}$ toàn cục thông qua tích cup, theo
đẳng cấu (1.5).

\fsection{3}{Các kết quả về đối đồng điều địa phương}

Cho $A$ là một vành giao hoán có đơn vị, cùng với một iđêan $J$. Với
$A$-môđun biến thiên $M$, ta sẽ định nghĩa các đồng cấu có tính hàm tử
\begin{equation}\tag{3.1}
\left\{
\begin{aligned}
\operatorname{Ext}_{A}^{p}(A/J,M)
  &\longrightarrow
  \operatorname{Hom}_{A}\!\left(\bigwedge^{p}J/J^{2},M\otimes A/J\right),\\
\operatorname{Tor}_{p}^{A}(A/J,M)
  &\longleftarrow
  \left(\bigwedge^{p}J/J^{2}\right)
  \otimes\operatorname{Hom}_{A}(A/J,M),
\end{aligned}
\right.
\end{equation}
trong đó các tích tenxơ và tích ngoài đều được lấy trên vành $A$; cũng lưu ý
rằng $J/J^{2}$ thực ra là một $A/J$-môđun, và các lũy thừa
ngoài của nó, xét như $A$-môđun hay như $A/J$-môđun, là như nhau.
Việc định nghĩa các đồng cấu (3.1) tương đương với việc định nghĩa, cho mọi
hệ $\underline{x}=(x_1,\ldots,x_p)$ gồm các phần tử của $J$, các đồng cấu
$\varphi_{\underline{x}}$:
\begin{equation}\tag{3.2}
\left\{
\begin{aligned}
\varphi_{\underline{x}}:
  \operatorname{Ext}_{A}^{p}(A/J,M)&\longrightarrow M\otimes A/J,\\
\varphi_{\underline{x}}:
  \operatorname{Hom}_{A}(A/J,M)&\longrightarrow
  \operatorname{Tor}_{p}^{A}(A/J,M),
\end{aligned}
\right.
\end{equation}
sao cho các điều kiện sau được thỏa mãn:


\src{174}
\begin{enumerate}
\item[\textup{i.}] $\varphi_{x_1,\ldots,x_p}$ phụ thuộc theo cách
$A$-đa tuyến tính thay phiên vào hệ các $x_i\in J$.
\item[\textup{ii.}] $\varphi_{x_1,\ldots,x_p}$ bằng không khi một trong các $x_i$ thuộc
$J^{2}$.
\end{enumerate}
Thực ra, ii. suy ra từ i., vì $a\varphi_{\underline{x}}=0$ với $a\in J$,
như thấy được khi lưu ý rằng mọi môđun trong (3.2) đều bị
$J$ triệt tiêu.

Để định nghĩa $\varphi_{\underline{x}}$, ta xét phức
$K_{\underline{x}}$ có $A$-môđun nền là $\bigwedge A^{p}$, và
có toán tử vi phân là phép co $i_{\underline{x}}$ theo
$\underline{x}$ được xem như một dạng tuyến tính trên $A^{p}$ có các thành phần
$x_1,\ldots,x_p$. Toán tử vi phân có bậc $-1$, các bậc đều
không âm, và nhóm đối đồng điều của phức này ở bậc $0$ là
$A/(x_1A+\cdots+x_pA)$. Vì các $x_i$ thuộc $J$, từ đó suy ra một
đồng cấu tăng cường $K_{\underline{x},0}\longrightarrow A/J$. Do đó,
$K_{\underline{x}}$ xuất hiện như một phức tự do tăng cường, với môđun
tăng cường $A/J$. Theo cách quen thuộc, từ đó suy ra các đồng cấu
\[
\operatorname{Ext}_{A}^{*}\bigl(H_0(K_{\underline{x}}),M\bigr)
\longrightarrow
H^{*}\bigl(\operatorname{Hom}_{A}(K_{\underline{x}},M)\bigr)
\]
và
\[
\operatorname{Tor}_{*}^{A}\bigl(H_0(K_{\underline{x}}),M\bigr)
\longleftarrow
H_{*}\bigl(K_{\underline{x}}\otimes_A M\bigr),
\]
do đó, bằng cách hợp thành với các đồng cấu trên các $\operatorname{Ext}$ và các
$\operatorname{Tor}$ suy ra từ đồng cấu tăng cường
$H_0(K_{\underline{x}})\longrightarrow A/J$, ta thu được các đồng cấu
\begin{equation}\tag{3.3}
\left\{
\begin{aligned}
\psi_{\underline{x}}:
\operatorname{Ext}_{A}^{*}(A/J,M)
&\longrightarrow H^{*}\bigl(\operatorname{Hom}_{A}(K_{\underline{x}},M)\bigr),\\
\psi_{\underline{x}}:
\operatorname{Tor}_{*}^{A}(A/J,M)
&\longleftarrow H_{*}\bigl(K_{\underline{x}}\otimes_A M\bigr).
\end{aligned}
\right.
\end{equation}
Ta kiểm tra ngay rằng, ở bậc cao nhất $p$, các nhóm (đối) đồng điều ở các vế
phải lần lượt là $M/(x_1M+\cdots+x_pM)$ (tương ứng, là tập các phần tử của $M$
bị từng $x_i$ triệt tiêu). Vì các $x_i$ thuộc $J$, từ đó suy ra các
đồng cấu
\begin{equation}\tag{3.4}
\left\{
\begin{aligned}
H^{p}\bigl(\operatorname{Hom}_{A}(K_{\underline{x}},M)\bigr)
&\longrightarrow M\otimes A/J,\\
H_{p}\bigl(K_{\underline{x}}\otimes_A M\bigr)
&\longleftarrow\operatorname{Hom}_{A}(A/J,M).
\end{aligned}
\right.
\end{equation}
Bằng cách hợp thành các đồng cấu (3.3) và (3.4), ta thu được các đồng cấu (3.2)
mà ta muốn định nghĩa. Việc kiểm tra i. khá dài dòng, nhưng không
gặp khó khăn nào.


\src{175}
\result{MỆNH ĐỀ}{4} Cho $A$ là một vành giao hoán có đơn vị, và cho
$(x_1,\ldots,x_p)$ là một dãy các phần tử của $A$ sao cho, với $1\leq i\leq p$,
ảnh của $x_i$ trong vành thương của $A$ theo iđêan sinh bởi
$(x_1,\ldots,x_{i-1})$ không phải là ước của $0$. Gọi $J$ là iđêan sinh
bởi các $x_i$. Khi đó $J/J^{2}$ là một $(A/J)$-môđun tự do có cơ sở là các
ảnh chính tắc của các $x_i$, phức $K_{\underline{x}}$ là một
phân giải tự do của $A/J$, và với mọi $A$-môđun $M$, các đồng cấu (3.1) ở
bậc $p$ đều song ánh. Điều tương tự cũng đúng với các đồng cấu tương tự
được định nghĩa ở mọi bậc $i$, miễn là $J.M=0$.

(Điểm cốt yếu, mà từ đó mọi khẳng định khác đều suy ra, là việc
$K_{\underline{x}}$ không có đồng điều ở các bậc dương, một sự kiện quen thuộc dưới các giả thiết đã nêu.)

\result{HỆ QUẢ}{1} Với $A$ và $J$ như trên, giả sử thêm rằng
$A$ là một đại số afin chính quy có $\dim n$ trên một trường hoàn hảo $k$, và
$A/J$ là một đại số afin chính quy. Ký hiệu $\Omega^{1}(A)$,
$\Omega^{1}(A/J)$ là các môđun vi phân Kähler. Khi đó có một
đẳng cấu chính tắc tương thích với phép địa phương hóa,
\begin{equation}\tag{3.5}
  \operatorname{Ext}_{A}^{p}\bigl(\Omega^{n-p}(A/J),\Omega^{n}(A)\bigr)=A/J.
\end{equation}
Thật vậy, $\Omega^{n-p}(A/J)$ là một $(A/J)$-môđun tự do hạng $1$, và tương tự,
$\Omega^{n}(A)$ là một $A$-môđun tự do hạng $n$,
do đó vế trái bằng
\[
  \operatorname{Ext}_{A}^{p}(A/J,A)
  \otimes\Omega^{n-p}(A/J)'\otimes\Omega^{n}(A)
\]
(trong đó dấu « $'$ » chỉ $(A/J)$-môđun đối ngẫu). Tích tenxơ của
hai thừa số cuối được đồng nhất với $\bigwedge^{p}(J/J^{2})$,
nên toàn bộ biểu thức
được đồng nhất với
\[
  \operatorname{Ext}_{A}^{p}\!\left(A/J,\bigwedge^{p}(J/J^{2})\right),
\]
và do đó, theo mệnh đề, với
\[
  \operatorname{Hom}_{A}\!\left(\bigwedge^{p}J/J^{2},
  \bigwedge^{p}J/J^{2}\right),
\]
tức là với $A/J$. Đặc biệt,
$\operatorname{Ext}_{A}^{p}\bigl(\Omega^{n-p}(A/J),\Omega^{n}(A)\bigr)$ chứa
một phần tử phân biệt tương ứng với phần tử đơn vị của $A/J$, được gọi là
\emph{lớp cơ bản của iđêan $J$ trong $A$}. (Thật ra, nó có thể được
định nghĩa dưới những điều kiện rộng hơn đáng kể.) Hệ quả 1 có thể được viết
dưới một dạng hình học hơn và toàn cục hơn:

\result{HỆ QUẢ}{2} Cho $X$ là một đa tạp không kỳ dị trên một trường đóng
đại số $k$, cho $Y$ là một đa tạp con đóng không kỳ dị của $X$,
$\Osh_X$ là bó các vành địa phương của $X$, và $\Osh_Y$ là bó các vành
địa phương của $Y$, được xem như một bó thương của $\Osh_X$. Cho $n$ là
số chiều của $X$ và $n-p$ là số chiều của $Y$.

\src{176}
Cho $\Omega_X$ (tương ứng, $\Omega_Y$) là bó các mầm của các dạng
vi phân chính quy trên $X$ (tương ứng, trên $Y$). Khi đó ta có các đẳng cấu
chính tắc:
\begin{equation}\tag{3.6}
  \Extsh_{\Osh_X}^{p}\bigl(\Omega_Y^{n-p},\Omega_X^{n}\bigr)=\Osh_Y,
\end{equation}
hoặc tương đương
\begin{equation}\tag{3.6 bis}
  \Extsh_{\Osh_X}^{p}\bigl(\Osh_Y,\Omega_X^{n}\bigr)=\Omega_Y^{n-p}.
\end{equation}
Công thức (3.6 bis) có thể dùng làm định nghĩa của $\Omega_Y^{n-p}$ khi $Y$
là một đa tạp kỳ dị. Chính xác hơn:

\result{MỆNH ĐỀ}{5} Cho $Y$ là một tập con đại số có chiều
$q=n-p$ trong một đa tạp đại số không kỳ dị $X$ có chiều $n$. Cho
$F$ là một bó đại số nhất quán trên $X$ có giá chứa trong $Y$, và $L$
là một bó đại số tự do địa phương trên $X$. Khi đó các bó
$\Extsh_{\Osh_X}^{i}(F,L)$ bằng không khi $i<p$, còn khi $i=p$ ta có một
đẳng cấu chính tắc
\begin{equation}\tag{3.7}
  \Extsh_{\Osh_X}^{p}(F,L)
  =\Homsh_{\Osh_X}\!\left(F,
  \Extsh_{\Osh_X}^{p}(\Osh_X/J,L)\right),
\end{equation}
trong đó $J$ là một bó iđêan tùy ý trên $X$ triệt tiêu $F$ và có
tập các không điểm là $Y$. Đặc biệt, nếu $F$ là một bó đại số
nhất quán trên $Y$, thì
\begin{equation}\tag{3.7 bis}
  \Extsh_{\Osh_X}^{p}(F,L)
  =\Homsh_{\Osh_Y}\!\left(F,\Extsh_{\Osh_X}^{p}(\Osh_Y,L)\right).
\end{equation}
Sau cùng, vẫn giả sử $F$ là một bó đại số nhất quán trên $Y$, các bó
\[
  E^{i}(F)=\Extsh_{\Osh_X}^{p+i}(F,\Omega_X^{n})
\]
không phụ thuộc vào cách nhúng không gian đại số $Y$ vào một đa tạp
không kỳ dị $X$.

Vì bài toán có tính địa phương, có thể giả sử $X$ là afin và $L=\Osh_X$. Điều này quy
bài toán về đại số giao hoán, thậm chí về đại số địa phương: nếu $A$ là một
vành địa phương chính quy và $M$ là một $A$-môđun mà giá của nó có chiều
$\leq q=n-p$, cần chứng minh rằng $\operatorname{Ext}_{A}^{i}(M,A)=0$ khi $i<p$, và
\[
  \operatorname{Ext}_{A}^{p}(M,A)
  =\operatorname{Hom}_{A}\!\left(M,\operatorname{Ext}_{A}^{p}(A/J,A)\right),
\]
trong đó $J$ là một iđêan tùy ý có « số chiều » $\leq q$ và triệt tiêu $M$. Đối với
khẳng định thứ nhất, ta quy nạp theo $q$: một phép rút gọn từng lớp trực tiếp đưa
về trường hợp $M$ có dạng $A/J$; sau đó, thay $J$ bằng một iđêan nhỏ hơn
và dùng giả thiết quy nạp cùng dãy khớp của các
$\operatorname{Ext}$, ta quy về trường hợp $J$ được sinh bởi một « hệ tham số »
như trong Mệnh đề 4, và khi đó kết quả là hiển nhiên. Kết quả
trên cho thấy rằng, nếu $J$ là một iđêan


\src{177}
cố định có « số chiều » $\leq q$, thì hàm tử phản biến
$E(M)=\operatorname{Ext}_{A}^{p}(M,A)$ trên phạm trù các $A/J$-môđun là
khớp trái; hơn nữa, nó biến các tổng trực tiếp thành các tích trực tiếp,
từ đó dễ dàng suy ra rằng
\[
  E(M)=\operatorname{Hom}_{A}\bigl(M,E(A)\bigr).
\]
 Sau cùng,
khẳng định cuối của Mệnh đề 5 tinh tế hơn và suy ra từ một
đặc trưng nội tại của các $E^{i}(F)$ nhờ một định lý đối ngẫu
địa phương mà không thể trình bày ở đây.

\result{HỆ QUẢ}{} Ký hiệu $\omega_Y^{q}$ là bó
$\Extsh_{\Osh_X}^{p}(\Osh_Y,\Omega_X^{n})$. Khi đó, đối với các bó đại số
nhất quán $F$ trên $Y$, ta có một đẳng cấu có tính hàm tử:
\begin{equation}\tag{3.8}
  \Extsh_{\Osh_X}^{p}(F,\Omega_X^{n})
  =\Homsh_{\Osh_X}(F,\omega_Y^{q}).
\end{equation}

\fsection{4}{Lớp đối đồng điều gắn với một đa tạp con}

Trong toàn bộ phần sau, $X$ ký hiệu một tập đại số có chiều $n$, được xác định
trên một trường $k$, mà để đơn giản ta giả sử là đóng đại số. Ngoại trừ
ở n\textsuperscript{o} 6, $X$ được giả sử không kỳ dị. Ta ký hiệu
$\Osh_X$ là bó các vành địa phương trên $X$, và
\[
  \Omega_X^{*}=\bigcup_p\Omega_X^{p}
\]
là bó các mầm của các dạng vi phân trên $X$. Nếu $Y$ là một tập con đóng của $X$, ta đồng nhất
các bó đại số nhất quán trên $Y$ với các bó nhất quán trên $X$
bằng không ngoài $Y$; điều này đặc biệt áp dụng cho $\Osh_Y$ và $\Omega_Y$.

\result{BỔ ĐỀ}{1} Cho $F$ là một bó đại số nhất quán trên $X$ mà
giá của nó có chiều $\leq n-p$, và cho $L$ là một bó đại số nhất quán
tự do địa phương trên $X$. Khi đó $\operatorname{Ext}_{\Osh_X}^{i}(X;F,L)$ bằng không
với $i<p$, và ta có một đẳng cấu chính tắc
\begin{equation}\tag{4.1}
  \operatorname{Ext}_{\Osh_X}^{p}(X;F,L)
  =H^{0}\bigl(X,\Extsh_{\Osh_X}^{p}(F,L)\bigr).
\end{equation}
Nếu $F$ là một bó đại số nhất quán trên một tập con đóng $Y$ của $X$ có
chiều $\leq n-p$, thì ta có một đẳng cấu chính tắc
\begin{equation}\tag{4.1 bis}
  \Extsh_{\Osh_X}^{p}(F,L)
  =\Homsh_{\Osh_X}\!\left(
    F\otimes L'\otimes\Omega_X^{n},\omega_Y^{n-p}
  \right),
\end{equation}
trong đó $\omega_Y^{n-p}$ là bó trên $Y$ được định nghĩa trong hệ quả của
Mệnh đề 5 (và đồng nhất với $\Omega_Y^{n-p}$ nếu $Y$ không kỳ dị).

Công thức (4.1) là hệ quả trực tiếp của dãy phổ trong
Mệnh đề 1 và của Mệnh đề 5; theo công thức (3.8), ta có thể
viết


\src{178}
\begin{align*}
\Extsh_{\Osh_X}^{p}(F,L)
&=L\otimes(\Omega_X^{n})'\otimes
  \Extsh_{\Osh_X}^{p}(F,\Omega_X^{n})\\
&=L\otimes(\Omega_X^{n})'\otimes
  \Homsh_{\Osh_X}(F,\omega_Y^{q})\\
&=\Homsh_{\Osh_X}\!\left(
  F\otimes L'\otimes\Omega_X^{n},\omega_Y^{q}\right),
\end{align*}
trong đó $q=n-p$, từ đó suy ra công thức (4.1 bis).

Đặc biệt, đặt $F=\Osh_Y$, $L=\Omega_X^{p}$; vì
$\Omega_X^{n}\otimes(\Omega_X^{p})'=\Omega_X^{n-p}$, ta thu được một đẳng cấu
chính tắc
\begin{equation}\tag{4.2}
  \operatorname{Ext}_{\Osh_X}^{p}(X;\Osh_Y,\Omega_X^{p})
  =\operatorname{Hom}_{\Osh_X}
  (X;\Omega_X^{n-p},\omega_Y^{n-p}).
\end{equation}
Bây giờ, để đơn giản, giả sử $Y$ không kỳ dị, do đó
$\omega_Y^{n-p}=\Omega_Y^{n-p}$. Ta có một đồng cấu tự nhiên từ
$\Omega_X^{n-p}$ đến $\Omega_Y^{n-p}$, từ đó suy ra một tiết diện chính tắc $s_Y$ của
bó $\Extsh_{\Osh_X}^{p}(\Osh_Y,\Omega_X^{p})$, và nếu
mọi thành phần của $Y$ đều có chiều $n-p$, ta gọi nó là \emph{tiết diện
cơ bản} của bó $\Extsh_{\Osh_X}^{p}(\Osh_Y,\Omega_X^{p})$. Theo (4.1), tiết diện
này xác định một phần tử của
$\operatorname{Ext}_{\Osh_X}^{p}(X;\Osh_Y,\Omega_X^{p})$. Mặt khác, đồng cấu
tự nhiên $\Osh_X\longrightarrow\Osh_Y$ xác định một đồng cấu
\[
  \operatorname{Ext}_{\Osh_X}^{p}(X;\Osh_Y,\Omega_X^{p})
  \longrightarrow
  \operatorname{Ext}_{\Osh_X}^{p}(X;\Osh_X,\Omega_X^{p})
  =H^{p}(X,\Omega_X^{p}).
\]
Như vậy, ta thu được một phần tử của $H^{p}(X,\Omega_X^{p})$, ký hiệu là $P_X(Y)$ và gọi là
\emph{lớp đối đồng điều của $Y$ trong $X$}. Do đó, lớp này được suy ra từ tiết diện
$s_Y$ của $\Extsh_{\Osh_X}^{p}(\Osh_Y,\Omega_X^{p})$ thông qua sơ đồ
các đồng cấu sau đây:
\begin{equation}\tag{4.3}
\begin{aligned}
H^{p}(X,\Omega_X^{p})
&=\operatorname{Ext}_{\Osh_X}^{p}(X;\Osh_X,\Omega_X^{p})\\
&\longleftarrow\operatorname{Ext}_{\Osh_X}^{p}(X;\Osh_Y,\Omega_X^{p})
\xrightarrow{\ \sim\ }
H^{0}\!\left(X,\Extsh_{\Osh_X}^{p}(\Osh_Y,\Omega_X^{p})\right).
\end{aligned}
\end{equation}

Ta gọi một \emph{chu trình không kỳ dị chiều $n-p$} là một phần tử của nhóm
Abel tự do sinh bởi các đa tạp con bất khả quy không kỳ dị có
chiều $n-p$ trong $X$. Khi đó, ánh xạ $Y\longmapsto P(Y)$ mở rộng thành
một đồng cấu từ nhóm các chu trình không kỳ dị chiều $n-p$ trên $X$
vào nhóm $H^{p}(X,\Omega_X^{p})$.

Cho $Z^{n-p}$ và $Z'^{n-p'}$ là các chu trình không kỳ dị lần lượt có chiều $n-p$ và
$n-p'$; ta nói chúng \emph{giao hoành} nếu mọi thành phần
của $Z$ giao hoành với mọi thành phần của $Z'$. Khi đó, chu trình $Z.Z'$ được
xác định và là một chu trình không kỳ dị chiều $n-p-p'$. Với quy ước đó, ta có


\src{179}
\result{ĐỊNH LÝ}{1} Nếu $Z^{n-p}$ và $Z'^{n-p'}$ là các chu trình không kỳ dị giao hoành với nhau, thì ta có
\begin{equation}\tag{4.4}
  P_X(Z.Z')=P_X(Z).P_X(Z'),
\end{equation}
trong đó tích ở vế phải là tích cup
\[
  H^{p}(X,\Omega_X^{p})\times H^{p'}(X,\Omega_X^{p'})
  \longrightarrow H^{p+p'}(X,\Omega_X^{p+p'}).
\]
(Ta giả sử $X$ đẳng cấu với một tập con đóng địa phương của một không gian
xạ ảnh.)

Giả thiết cuối này chỉ được dùng để suy ra rằng mọi bó đại số nhất quán trên $X$
là thương của một bó đại số nhất quán tự do địa phương (SERRE),
và do đó có một phân giải trái bởi các bó tự do địa phương. Để chứng minh
Định lý 1, ta có thể giả sử $Z$ và $Z'$ là các đa tạp con bất khả quy không kỳ dị
$Y$ và $Y'$ giao hoành với nhau. Cho $L_{*}$ là một phân giải trái
của $\Osh_Y$ bởi
các bó tự do địa phương; khi đó, theo Mệnh đề 3, biểu đồ các đồng cấu
(4.3) đồng nhất với biểu đồ
\[
R^{p}\Gamma(\Omega_X^{p})
\xleftarrow{\alpha}
\mathbb R^{p}\Gamma\!\left(\Homsh_{\Osh_X}(L_{*},\Omega_X^{p})\right)
\xrightarrow[\sim]{\beta}
\Gamma\!\left(H^{p}\bigl(\Homsh_{\Osh_X}(L_{*},\Omega_X^{p})\bigr)\right),
\]
trong đó $\Gamma$ là hàm tử “nhóm các tiết diện” trên phạm trù
các bó nhóm Abel trên $X$, còn $\mathbb R^{p}\Gamma$ chỉ
siêu đối đồng điều bậc $p$ của hàm tử ấy, và $R^{p}\Gamma$ là
hàm tử dẫn xuất thứ $p$ của nó. Để đơn giản, ta giả sử $L_0=\Osh_X$ và
đồng cấu tăng cường $L_0\longrightarrow\Osh_Y$ là
đồng cấu tự nhiên (điều này được phép); khi đó $\alpha$ được suy ra
từ đồng cấu
các phức $\Osh_X\longrightarrow L_{*}$ (trong đó $\Osh_X$ được xem như một phức
tập trung ở bậc $0$), có tính đến rằng $\mathbb R^{p}\Gamma(K)=R^{p}\Gamma(K_0)$ nếu $K$
là một phức các bó tập trung ở bậc $0$. Đồng cấu $\beta$ là một đồng cấu
biên quen thuộc. Xét một biểu đồ tương tự, ứng với
một phân giải tự do địa phương
$L'_{*}$ của $\Osh_{Y'}$, và xét biểu đồ giao hoán của các phép ghép cặp:
\begin{equation}\tag{4.5}
\resizebox{\textwidth}{!}{$
\begin{array}{ccc}
R^{p}\Gamma(\Omega_X^{p})
&\longleftarrow
\mathbb R^{p}\Gamma\!\left(\Homsh_{\Osh_X}(L_{*},\Omega_X^{p})\right)
&\xrightarrow{\sim}
\Gamma H^{p}\!\left(\Homsh_{\Osh_X}(L_{*},\Omega_X^{p})\right)
\\[3pt]
\times&\times&\times\\[3pt]
R^{p'}\Gamma(\Omega_X^{p'})
&\longleftarrow
\mathbb R^{p'}\Gamma\!\left(\Homsh_{\Osh_X}(L'_{*},\Omega_X^{p'})\right)
&\xrightarrow{\sim}
\Gamma H^{p'}\!\left(\Homsh_{\Osh_X}(L'_{*},\Omega_X^{p'})\right)
\\[3pt]
\downarrow&\downarrow&\downarrow\\[3pt]
R^{p+p'}\Gamma(\Omega_X^{p+p'})
&\longleftarrow
\mathbb R^{p+p'}\Gamma\!\left(
\Homsh_{\Osh_X}(L_{*}\otimes L'_{*},\Omega_X^{p+p'})\right)
&\xrightarrow{\sim}
\Gamma H^{p+p'}\!\left(
\Homsh_{\Osh_X}(L_{*}\otimes L'_{*},\Omega_X^{p+p'})\right).
\end{array}
$}
\end{equation}


\src{180}
Các phép ghép cặp của hai cột bên phải được suy ra từ phép ghép cặp của
các phức bó
\[
\Homsh_{\Osh_X}(L_{*},\Omega_X^{p})
\times\Homsh_{\Osh_X}(L'_{*},\Omega_X^{p'})
\longrightarrow
\Homsh_{\Osh_X}(L_{*}\otimes L'_{*},\Omega_X^{p+p'}),
\]
được định nghĩa bằng tích ngoài
$\Omega_X^{p}\times\Omega_X^{p'}\longrightarrow\Omega_X^{p+p'}$.
Phép ghép cặp của cột thứ nhất là tích cup (tương ứng với tích
ngoài). Tôi khẳng định hàng cuối của (4.5) đồng nhất với biểu đồ
các đẳng cấu tương tự (4.3), trong đó $Y$ được thay bởi $Y\cap Y'$ và $p$ bởi
$p+p'$. Để làm điều này, chỉ cần chứng minh rằng $L\otimes L'$ là một phân giải
(hiển nhiên tự do địa phương) của $\Osh_{Y\cap Y'}$. Thật vậy, ta có
\[
  H_0(L\otimes L')=\Osh_Y\otimes\Osh_{Y'}=\Osh_{Y\cap Y'}
\]
và
\[
  H_i(L\otimes L')=\operatorname{Tor}_{i}^{\Osh_X}(\Osh_Y,\Osh_{Y'})=0
  \qquad(i>0),
\]
vì $Y$ và $Y'$ giao hoành với nhau. Định lý 1 giờ đây suy ra
từ công thức:
\begin{equation}\tag{4.6}
  s_Y.s_{Y'}=s_{Y.Y'}
\end{equation}
(trong đó tích ở vế trái là tích của cột cuối trong (4.5)).
Công thức (4.6) này, có bản chất thuần túy địa phương, được kiểm tra dễ dàng
bằng cách lấy $L_{*}$ và $L'_{*}$ là các phân giải đã xét trong Mệnh đề
4. Tương tự, ta chứng minh (bằng một chứng minh đơn giản hơn) rằng
$Z\longmapsto P_X(Z)$ tương thích với tích Đề-các:
\begin{equation}\tag{4.7}
  P_{X\times X'}(Z\times Z')=P_X(Z)\otimes P_{X'}(Z')
\end{equation}
(công thức có hiệu lực nếu $Z$, tương ứng $Z'$, là một chu trình không kỳ dị trên đa tạp
không kỳ dị $X$, tương ứng $X'$, và $Z\times Z'$ được xem như một chu trình không
kỳ dị trên $X\times X'$). Từ (4.4) và (4.7), suy ra rằng $P_X(Z)$ cũng
tương thích với phép “lấy ảnh ngược” theo một cấu xạ của các đa tạp không
kỳ dị $f:X\longrightarrow X'$:
\begin{equation}\tag{4.8}
  P_X\bigl(f^{-1}(Z')\bigr)=f^{*}\bigl(P_{X'}(Z')\bigr),
\end{equation}
công thức có hiệu lực nếu $Z'$ là một chu trình không kỳ dị trên $X'$ sao cho $f$
“giao hoành” với $Z'$, nghĩa là đồ thị của $f$ giao hoành với chu trình
$X\times Z'$ trong $X\times X'$.


\src{181}
\result{HỆ QUẢ}{1} Cho $X$, $X'$ là hai đa tạp không kỳ dị
đóng địa phương trong một không gian xạ ảnh, và giả sử $X'$ đầy đủ. Cho $U$
là một chu trình không kỳ dị trên $X\times X'$, và $a$, $b$ là hai điểm của $X'$ sao cho
$U$ giao hoành với các chu trình $X\times(a)$ và $X\times(b)$; cho $Z$
và $Z'$ là các chu trình không kỳ dị trên $X$ sao cho
\[
  Z\times(a)=(X\times(a)).U,
  \qquad
  Z'\times(b)=(X\times(b)).u.
\]
Khi đó ta có
\[
  P_X(Z)=P_X(Z').
\]

Thật vậy, cho $f_a:X\longrightarrow X\times X'$ được xác định bởi
$f_a(x)=(x,a)$; khi đó, theo (4.8), ta có công thức
$P(Z)=f_a^{*}(P(U))$, và tương tự $P(Z')=f_b^{*}(P(U))$. Mặt khác, dùng công thức
Künneth
\[
  H^{*}(X\times X',\Omega_{X\times X'}^{*})
  =H^{*}(X,\Omega_X^{*})\otimes H^{*}(X',\Omega_{X'}^{*})
\]
và sự kiện $H^{0}(X',\Osh_{X'})$ chỉ gồm các vô hướng, ta dễ dàng thu được
$f_a^{*}=f_b^{*}$, từ đó suy ra kết luận.

Với mọi $x\in X$, $(x)$ là một đa tạp con không kỳ dị của $X$ có
đối chiều $n$, và do đó xác định một phần tử $\varepsilon_x$ của
$H^{n}(X,\Omega_X^{n})$. Nếu $X$ là một đa tạp xạ ảnh không kỳ dị, từ
Hệ quả 1 suy ra rằng $\varepsilon_x$ không phụ thuộc vào điểm $x$ được chọn; ta ký hiệu nó
bởi $\varepsilon_X$ và gọi là \emph{lớp cơ bản} của
$H^{n}(X,\Omega_X^{n})$.

\result{NHẬN XÉT}{} Để có một lý thuyết thỏa đáng, cần định nghĩa
$P_X(Z)$ cho các chu trình $Z$ tùy ý và chứng minh Định lý 1 cho một
giao đúng của các chu trình. (Vào thời điểm viết bài trình bày này, điều đó vẫn chưa được
thực hiện trong tính tổng quát hoàn toàn). Giả sử điều đó đã được thực hiện, Hệ quả 1
trở thành: nếu $Z$ và $Z'$ là hai chu trình tương đương đại số, thì
$P_X(Z)=P_X(Z')$ (một mệnh đề dường như không suy ra từ những điều trên, ngay cả khi
$Z$ và $Z'$ không kỳ dị).

\fsection{5}{Định lý đối ngẫu}

Trong mục này, $X$ ký hiệu một đa tạp xạ ảnh không kỳ dị có chiều
$n$.

\result{ĐỊNH LÝ}{2} Lớp cơ bản $\varepsilon_X$ của
$H^{n}(X,\Omega_X^{n})$ là một cơ sở của không gian vectơ này.

(Chứng minh sẽ được trình bày ở dưới). Nhờ định lý trên, ta có thể
đồng nhất $H^{n}(X,\Omega_X^{n})$ với trường $k$. Bây giờ xét các
phép ghép cặp đã xét ở n\textsuperscript{o} 2, chúng đặc biệt cho một
phép ghép cặp


\src{182}
\[
\operatorname{Ext}_{\Osh_X}^{p}(X;\Osh_X,F)
\times\operatorname{Ext}_{\Osh_X}^{n-p}(X;F,\Omega_X^{n})
\longrightarrow
\operatorname{Ext}_{\Osh_X}^{n}(X;\Osh_X,\Omega_X^{n}),
\]
tức là,
\begin{equation}\tag{5.1}
  H^{p}(X,F)
  \times\operatorname{Ext}_{\Osh_X}^{n-p}(X;F,\Omega_X^{n})
  \longrightarrow H^{n}(X,\Omega_X^{n})
\end{equation}
Theo Định lý 2, phép ghép cặp này xác định một đồng cấu
\begin{equation}\tag{5.2}
  \operatorname{Ext}_{\Osh_X}^{n-p}(X;F,\Omega_X^{n})
  \longrightarrow\bigl(H^{p}(X,F)\bigr)'
\end{equation}
Đồng cấu này có tính hàm tử theo $F$ và giao hoán với các đồng cấu nối
tương ứng với các dãy khớp $0\longrightarrow F'\longrightarrow F
\longrightarrow F''\longrightarrow0$.

\result{ĐỊNH LÝ}{3} Đồng cấu (5.2) là một đẳng cấu.

Đặc biệt, ta thu lại kết quả của Serre:

\result{HỆ QUẢ}{} Cho $E$ là một phân thớ vectơ đại số trên $X$ và
$\Osh_X(E)$ là bó các mầm của các tiết diện chính quy của $E$. Khi đó ta có các
đẳng cấu chính tắc:
\begin{equation}\tag{5.3}
  \bigl(H^{p}(X,\Osh_X(E))\bigr)'
  =H^{n-p}\bigl(X,\Omega_X^{n}\otimes\Osh_X(E')\bigr)
\end{equation}
Chỉ cần áp dụng Định lý 3 và Hệ quả 1 cho Mệnh đề 1.

Các Định lý 2 và 3 sẽ được suy ra từ mệnh đề sau:

\noindent\textbf{(D) Đồng cấu sau đây}
\begin{equation}\tag{5.2 bis}
  \operatorname{Ext}_{\Osh_X}^{n-p}(X;F,\Omega_X^{n})
  \longrightarrow\bigl(H^{p}(X,F)\bigr)'\otimes L,
  \qquad L=H^{n}(X,\Omega_X^{n}),
\end{equation}
được cảm sinh bởi phép ghép cặp (5.1) là một đẳng cấu.

Quả vậy, ta chứng minh rằng (D) kéo theo Định lý 2. Cho
$k_x=\Osh_{(x)}$ là bó các vành địa phương của đa tạp được thu về điểm
$x\in X$. Xét đồng cấu chính tắc
$\Osh_X\longrightarrow k_x$ và đồng cấu tương ứng
\begin{equation}\tag{5.3}
  H^{0}(X,\Osh_X)\longrightarrow H^{0}(X,k_x)
\end{equation}
chuyển vị của đồng cấu này được đồng nhất với đồng cấu
\begin{equation}\tag{5.4}
  \operatorname{Ext}_{\Osh_X}^{n}(X;k_x,\Omega_X^{n})\otimes L'
  \longrightarrow
  \operatorname{Ext}_{\Osh_X}^{n}(X;\Osh_X,\Omega_X^{n})\otimes L'
\end{equation}
được cảm sinh bởi đồng cấu trên các $\operatorname{Ext}^{n}$ tương ứng với
$\Osh_X\longrightarrow k_x$:
\begin{equation}\tag{5.5}
  \operatorname{Ext}_{\Osh_X}^{n}(X;k_x,\Omega_X^{n})
  \longrightarrow
  \operatorname{Ext}_{\Osh_X}^{n}(X;\Osh_X,\Omega_X^{n})
  =H^{n}(X,\Omega_X^{n})
\end{equation}


\src{183}
Vì (5.3) là một đẳng cấu, nên (5.4) cũng là một đẳng cấu, và do đó (5.5)
cũng vậy. Vì $s_{(x)}$ là một cơ sở của
$\operatorname{Ext}_{\Osh_X}^{n}(X;k_x,\Omega_X^{n})$ theo (4.2), suy ra ảnh
$\varepsilon_X$ của nó là một cơ sở của $H^{n}(X,\Omega_X^{n})$.

Còn phải chứng minh mệnh đề (D); mệnh đề này sẽ suy ra một cách hoàn toàn hình
thức từ các sự kiện sơ cấp được tóm tắt trong những bổ đề sau. Trong đó, ta giả
sử $X$ là một phần đóng (kỳ dị hoặc không) của không gian xạ ảnh $P$ chiều
$r$. Ta dùng ký hiệu $\Osh_P(m)$ cho bó trên $P$ được ký hiệu $\Osh(m)$ trong
[3], và ký hiệu $\Osh_X(m)$ cho bó tương tự trên $X$.

\result{BỔ ĐỀ}{2} Mệnh đề (D) đúng nếu $X=P$ và $F=\Osh_P(m)$.

Bổ đề này được kiểm tra bằng một phép tính trực tiếp. Phép tính tường minh các
$H^{i}(P,\Osh_P(m))$ có trong [3], nhưng cũng có thể thực hiện một cách sơ cấp
hơn. Phép tính tích cup
\[
  H^{i}(P,\Osh_P(m))\times H^{j}(P,\Osh_P(m'))
  \longrightarrow H^{i+j}(P,\Osh_P(r+r'))
\]
cần thiết để viết tường minh phép ghép cặp (5.1) không gây khó khăn.

\result{BỔ ĐỀ}{3} Mọi bó đại số nhất quán $F$ trên $X$ đều đẳng cấu với một
bó thương của một bó $\Osh_X(-m)^{k}$, trong đó có thể giả sử $m$ lớn tùy ý.

Điều này suy ra từ sự kiện rằng $F\otimes\Osh_X(m)$ ``được sinh bởi các tiết
diện'' khi $m$ lớn; xem [3].

\result{BỔ ĐỀ}{4} Cho $i>0$. Khi đó
$H^{r-i}(P,\Osh_P(-m))=0$ với $m$ lớn; và với mọi bó đại số nhất quán $B$
trên $X$, ta có
\[
  \operatorname{Ext}_{\Osh_X}^{i}(X;\Osh_X(-m),B)=0
\]
với $m$ lớn.

Khẳng định thứ nhất suy ra từ những phép tính tường minh đã nhắc ở trên; đối
với khẳng định thứ hai, ta nhận thấy có một đẳng cấu
\[
  \operatorname{Ext}_{\Osh_X}^{i}(X;\Osh_X(-m),B)
  =H^{i}(X,B\otimes\Osh_X(m))
\]
(Mệnh đề 1, Hệ quả 1), từ đó kết luận suy ra nhờ một kết quả quen thuộc trong
[3]. Kết hợp hai bổ đề trên, ta thu được

\result{HỆ QUẢ}{} Cho $i>0$. Khi đó hàm tử
$F\longmapsto H^{r-i}(P,F)$ trên phạm trù các bó đại số nhất quán trên $P$
là xóa được, và hàm tử
$\operatorname{Ext}_{\Osh_X}^{i}(X;F,B)$ trên phạm trù các bó đại số nhất quán
trên $X$ cũng vậy.


\src{184}
\result{BỔ ĐỀ}{5} Cho $A$ và $B$ là hai bó đại số nhất quán trên
$X$, và đặt $A(m)=A\otimes\Osh_X(m)$. Khi đó, với $m$ đủ lớn,
đồng cấu chính tắc
\[
  \operatorname{Ext}_{\Osh_X}^{1}(X;A(-m),B)
  \longrightarrow H^{0}\bigl(X,\Extsh_{\Osh_X}^{1}(A(-m),B)\bigr)
  =H^{0}\bigl(X,\Extsh_{\Osh_X}^{1}(A,B)(m)\bigr)
\]
là một đẳng cấu.

Điều này suy ra ngay từ dãy phổ của Mệnh đề 1 áp dụng cho
$A(-m)$ và $B$, vì khi đó ta có
\[
  E_{2}^{p,q}(A(-m),B)
  =H^{p}\bigl(X,\Extsh_{\Osh_X}^{q}(A(-m),B)\bigr)
  =H^{p}\bigl(X,\Extsh_{\Osh_X}^{q}(A,B)(m)\bigr),
\]
và biểu thức này bằng không khi $p>0$ và $m$ lớn.

Bây giờ ta chứng minh (D) trong trường hợp $X=P$. Trước hết, ta chứng minh (5.2 bis)
là một đẳng cấu khi $p=n$; vì khi đó cả hai vế đều là các
hàm tử khớp trái (do $H^{r+1}(P,F)=0$), từ Bổ đề 3 suy ra rằng
chỉ cần chứng minh khẳng định cho $F=\Osh_P(-m)$; nhưng trong trường hợp này, khẳng định đó
đã nằm trong Bổ đề 2. Vì các đồng cấu (5.2 bis) có tính hàm tử và
tương thích với các toán tử nối, và vì khi $p<n$, cả hai vế của
(5.2 bis), xét như các hàm tử theo $F$, đều xóa được (Hệ quả của Bổ đề 4), nên
một lập luận chuẩn cho thấy (5.2 bis) là một đẳng cấu
với mọi $p$. Điều này chứng minh định lý đối ngẫu cho không gian xạ ảnh.
Bây giờ giả sử $X$ tùy ý nhưng không kỳ dị. Theo định lý
đối ngẫu cho $P$, ta có một đẳng cấu
\[
  H^{n}(X,F)'=H^{n}(P,F)'
  =\operatorname{Ext}_{\Osh_P}^{r-n}(P;F,\Omega_P^{r}).
\]
Theo Bổ đề 1 (n\textsuperscript{o} 4), vế cuối được đồng nhất với
\[
  \operatorname{Hom}_{\Osh_P}(P;F,\omega_X^{n})
  =\operatorname{Hom}_{\Osh_X}(X;F,\Omega_X^{n})
  =\operatorname{Ext}_{\Osh_X}^{0}(X;F,\Omega_X^{n}),
\]
từ đó thu được một đẳng cấu
\begin{equation}\tag{5.6}
  H^{n}(X,F)'=\operatorname{Hom}_{\Osh_X}(X;F,\Omega_X^{n})
  =\operatorname{Ext}_{\Osh_X}^{0}(X;F,\Omega_X^{n}).
\end{equation}
Cho $F=\Omega_X^{n}$, ta thu được một đẳng cấu
\begin{equation}\tag{5.7}
  \gamma:H^{n}(X,\Omega_X^{n})\xrightarrow{\sim}k.
\end{equation}


\src{185}
Ta kiểm tra rằng đẳng cấu (5.6) chính là (5.2 bis) với $p=n$, khi
đặt $L=k$ nhờ (5.7). Do đó, (5.2 bis) là một đẳng cấu với
$p=n$. Để chứng minh đây là một đẳng cấu với mọi $p$, một lần nữa chỉ cần
chứng minh rằng, với $p<n$, hai vế của (5.2 bis) là các hàm tử
xóa được theo $F$; hơn nữa, theo Bổ đề 3, chỉ cần chứng minh mệnh đề mạnh hơn
rằng cả hai vế bằng không khi đặt $F=\Osh_X(-m)$ với $m$ lớn. Đối với vế
thứ nhất, điều này đúng theo Bổ đề 4; còn đối với vế thứ hai, dùng
định lý đối ngẫu trên $P$, ta viết:
\[
  H^{p}(X,\Osh_X(-m))'
  =\operatorname{Ext}_{\Osh_P}^{r-p}
    (P;\Osh_X(-m),\Omega_P^{r}).
\]
Vế thứ hai bằng không khi $p<n$ và $m$ lớn, như suy ra từ Bổ đề 5
(trong đó ta đặt $X=P$) và từ việc $\Osh_X$, xét như một bó đại số
nhất quán trên $P$, có chiều đối đồng điều $\leq r-n$ (vì $X$ không
kỳ dị); do đó
\[
  \Extsh_{\Osh_P}^{r-p}(\Osh_X,\Omega_P^{r})=0
  \qquad\text{với }p<n.
\]

\fsection{6}{Định lý đối ngẫu cho các đa tạp kỳ dị}

Cho $X$ là một tập con đóng có chiều $n$ của không gian xạ ảnh $P$ có
chiều $r$. Ở đây, công thức (5.6) phải được viết dưới dạng
\begin{equation}\tag{6.1}
  H^{n}(X,F)'
  \simeq\operatorname{Hom}_{\Osh_X}(X;F,\omega_X^{n})
  =\operatorname{Ext}_{\Osh_X}^{0}(X;F,\omega_X^{n}),
\end{equation}
trong đó ta đặt
\begin{equation}\tag{6.2}
  \omega_X^{n}=E^{0}(\Osh_X)
  =\Extsh_{\Osh_P}^{r-n}(\Osh_X,\Omega_P^{r}).
\end{equation}
Như đã nói trong Mệnh đề 5, bó vừa được đưa vào thực ra không phụ thuộc
vào phép nhúng đã chọn của $X$ vào một đa tạp không kỳ dị $P$.
Đặt $F=\omega_X^{n}$ trong (6.1), ta được
\begin{equation}\tag{6.2}
  H^{n}(X,\omega_X^{n})'
  \simeq\operatorname{Hom}_{\Osh_X}(X;\omega_X^{n},\omega_X^{n}),
\end{equation}
 từ đó tồn tại một phần tử
phân biệt trong $H^{n}(X,\omega_X^{n})'$, tương ứng với đồng cấu
đồng nhất từ $\omega_X^{n}$ vào chính nó:
\begin{equation}\tag{6.3}
  \gamma:H^{n}(X,\omega_X^{n})\longrightarrow k.
\end{equation}

\src{186}
Khi đó, ta xét các phép ghép cặp được xác định bởi phép hợp thành các
$\operatorname{Ext}$:
\begin{equation}\tag{6.4}
  H^{p}(X,F)\times
  \operatorname{Ext}_{\Osh_X}^{n-p}(X;F,\omega_X^{n})
  \longrightarrow H^{n}(X,\omega_X^{n}),
\end{equation}
rồi hợp thành chúng với đồng cấu $\gamma$ trong (6.3); ta thu được
các đồng cấu có tính hàm tử, tương thích với các toán tử nối:
\begin{equation}\tag{6.5}
  \operatorname{Ext}_{\Osh_X}^{n-p}(X;F,\omega_X^{n})
  \longrightarrow H^{p}(X,F)'
\end{equation}
(tổng quát hóa (5.2)). Ta kiểm tra rằng khi $p=n$, cách này cho
đẳng cấu (6.1). Như vậy, ta có

\result{ĐỊNH LÝ}{3 bis} Bốn điều kiện sau đây về $X$ là tương đương đối với
một số nguyên $k\geq0$ cho trước:
\begin{enumerate}
  \item[\textit{i.}] Đồng cấu có tính hàm tử (6.5) là một đẳng cấu với
    $n-k\leq p\leq n$.
  \item[\textit{ii.}] Ta có $H^{p}(X,\Osh_X(-m))=0$ với $m$ lớn và với
    $n-k\leq p<n$.
  \item[\textit{iii.}] Hàm tử $H^{p}(X,F)$ trên phạm trù các bó đại số
    nhất quán trên $X$ là xóa được với $n-k\leq p<n$.
  \item[\textit{iv.}] Ta có
    $E^{i}(\Osh_X)=\Extsh_{\Osh_P}^{r-n+i}(\Osh_X,\Omega_P^{r})=0$
    với $0<i\leq k$.
\end{enumerate}

\noindent\textbf{CHỨNG MINH.} --- \textit{i.}$\Rightarrow$\textit{ii.} theo
Bổ đề 4, \textit{ii.}$\Rightarrow$\textit{iii.} theo Bổ đề 3, còn
\textit{iii.}$\Rightarrow$\textit{i.} bằng một lập luận chuẩn quen thuộc,
vì khi đó hai vế của (6.5), xét như các hàm tử, đều xóa được với
$n-k\leq p<n$ (vế thứ nhất là như vậy theo Bổ đề 4). Cuối cùng, ta chứng minh
\textit{ii.}$\Leftrightarrow$\textit{iv.} Điều này suy ra từ hệ quả của
mệnh đề sau đây:

\result{MỆNH ĐỀ}{6} Cho $F$ là một bó đại số nhất quán trên $X$, và cho
$i$ là một số nguyên. Khi đó, với $m$ đủ lớn, ta có một đẳng cấu:
\begin{equation}\tag{6.6}
  H^{i}(X,F(-m))'
  \simeq H^{0}\bigl(X,E^{n-i}(F)(m)\bigr),
\end{equation}
trong đó ta đặt (so sánh n\textsuperscript{o} 3, Mệnh đề 5):
\begin{equation}\tag{6.7}
  E^{j}(F)=\Extsh_{\Osh_P}^{r-n+j}(F,\Omega_P^{r}).
\end{equation}
Thật vậy, theo định lý đối ngẫu trên $P$, vế thứ nhất của (6.6)
đẳng cấu với
$\operatorname{Ext}_{\Osh_P}^{r-i}(P;F(-m),\Omega_P^{r})$, do đó (6.6) suy ra
từ Bổ đề 5.

\src{187}
\result{HỆ QUẢ}{} Ta có $H^{i}(X,F(-m))=0$ với $m$ đủ lớn nếu và chỉ nếu
$E^{n-i}(F)=0$.

Nhắc lại rằng các $E^{j}(F)$ thực ra không phụ thuộc vào phép nhúng xạ ảnh
đang xét. Điều kiện trong hệ quả có tính hoàn toàn địa phương; do đó,
nếu nó đúng đối với $F$, thì nó cũng đúng đối với mọi bó đẳng cấu địa phương
với một bó $F^{n}$. Đặc biệt, nếu điều kiện này đúng
đối với $\Osh_X$, thì nó đúng đối với mọi bó đại số nhất quán tự do địa phương.
Chẳng hạn, điều này đúng với mọi $i<n$ nếu $X$ không kỳ dị, với
$i=0$ nếu không thành phần nào của $X$ chỉ gồm một điểm duy nhất, và với $i=0,1$
nếu $S$ là chuẩn tắc và mọi thành phần của nó có chiều $>1$ (xem
[3]). Để điều này đúng với mọi
$i<k$, theo định nghĩa, cần và đủ là các vành địa phương
$\Osh_x$ ($x\in X$) có « đối chiều đồng điều » $\geq k$ (xem [4]
về khái niệm này). Nếu $k=n$, theo Định lý 3 bis điều đó có nghĩa rằng
định lý đối ngẫu đúng cho $X$, nghĩa là (6.5) là một đẳng cấu với
mọi $p$ và mọi $F$. Có thể đưa ra nhiều điều kiện tương đương trên
các vành địa phương $\Osh_x$ để điều này xảy ra (Nagata), chẳng hạn chúng
thỏa mãn định lý đẳng chiều Cohen--Macaulay. Điều này xảy ra,
chẳng hạn, nếu về địa phương $X$ là một « giao đầy đủ » trong một
đa tạp xung quanh không kỳ dị.

\fsection{7}{Đối ngẫu Poincaré}

Cho $X$ là một đa tạp xạ ảnh không kỳ dị có chiều $n$. Khi đó
\[
  H^{**}(X)=H^{*}(X,\Omega_X^{*})
\]
là một đại số phản giao hoán song phân bậc hữu hạn chiều, được phân bậc
theo bậc toàn phần, sao cho $H^{p,q}(X)=H^{p}(X,\Omega_X^{q})$ có bậc
$p+q$. Các bậc của $H^{*}(X)$ nằm giữa $0$ và $2n$. Theo
Định lý 2 và hệ quả của Định lý 3, $H^{*}(X)$ là một « đại số
Poincaré » có chiều $2n$, nghĩa là $H^{2n}(X)$ được trang bị một đẳng cấu với
trường cơ sở $k$, và tích
\[
  H^{m}(X)\times H^{2n-m}(X)\longrightarrow H^{2n}(X)=k
\]
là một phép đối ngẫu giữa $H^{m}(X)$ và $H^{2n-m}(X)$. Hơn nữa, nếu $Y$ là một
đa tạp xạ ảnh không kỳ dị khác, công thức Künneth cho các
bó đại số nhất quán cho
\begin{equation}\tag{7.1}
  H^{*}(X\times Y)=H^{*}(X)\otimes H^{*}(Y),
\end{equation}
một đẳng cấu tương thích với các cấu trúc đại số Poincaré. Hơn nữa,
$H^{*}(X)$, xét như một đại số giao hoán, là một hàm tử phản biến theo
$X$, vì một cấu xạ $f:Y\longrightarrow X$ hiển nhiên xác định một
đồng cấu của các đại số phân bậc

\src{188}
\begin{equation}\tag{7.2}
  f^{*}:H^{*}(X)\longrightarrow H^{*}(Y).
\end{equation}
Vì đây là các đại số Poincaré, phép chuyển vị cho ta một đồng cấu giữa các
không gian vectơ:
\begin{equation}\tag{7.3}
  f_{*}:H^{*}(Y)\longrightarrow H^{*}(X).
\end{equation}
Ta đã thấy ở n\textsuperscript{o} 4 rằng tác động của $f^{*}$ lên các lớp đối đồng điều
ứng với các chu trình không kỳ dị được diễn giải về mặt hình học bằng cách
lấy các lớp đối đồng điều ứng với ảnh ngược của chúng. Trong trường hợp hiện
tại, cần chứng minh rằng (7.3) cũng ứng với phép « ảnh trực tiếp » trên các
chu trình. Điều này suy ra, ít nhất dưới các điều kiện không kỳ dị thích hợp,
từ trường hợp riêng sau đây:

\result{ĐỊNH LÝ}{4} Nếu $f$ là ánh xạ đồng nhất của một đa tạp con không
kỳ dị $Y^{m}$ của $X^{n}$ vào $X^{n}$, thì, ký hiệu
phần tử đơn vị của $H(Y)$ là $1_Y$, ta có
\begin{equation}\tag{7.4}
  f_{*}(1_Y)=P_X(Y),
\end{equation}
trong đó vế thứ hai là lớp đối đồng điều trong $X$ gắn với $Y$.

Công thức này tương đương với
\begin{equation}\tag{7.4 bis}
  \bigl\langle\xi^{m,m},P_X(Y)\bigr\rangle\varepsilon_Y
  =f_{*}(\xi^{m,m})
  \qquad
  \bigl(\xi^{m,m}\in H^{m}(X,\Omega_X^{m})\bigr),
\end{equation}
 trong đó $\varepsilon_Y$ là phần tử cơ bản
trong $H^{m}(Y,\Omega_Y^{m})$. Trong trường hợp các đa tạp xạ ảnh không kỳ dị,
công thức này cho một định nghĩa mới về lớp đối đồng điều gắn với $Y$. Để chứng
minh Định lý 4, nhờ Định lý 3, ta diễn giải chuyển vị của đồng cấu
\[
  H^{m}(X,\Omega_X^{m})
  \longrightarrow H^{m}(Y,\Omega_Y^{m})
  =H^{m}(X,\Omega_Y^{m})
\]
như đồng cấu
\begin{equation}\tag{7.5}
\begin{split}
  H^{n-m}(X,\Omega_X^{n-m})
  &\simeq
  \operatorname{Ext}_{\Osh_X}^{n-m}(X;\Omega_X^{m},\Omega_X^{n})\\
  &\longleftarrow
  \operatorname{Ext}_{\Osh_X}^{n-m}(X;\Omega_Y^{m},\Omega_X^{n})
  \simeq
  \operatorname{Hom}_{\Osh_X}(X;\Omega_Y^{m},\Omega_Y^{m}).
\end{split}
\end{equation}
Ta kiểm tra rằng phần tử $1_Y$ của đối ngẫu của $H^{m}(Y,\Omega_Y^{m})$
được đồng nhất với phần tử ở vế phải tương ứng với tự đồng cấu đồng nhất của
$\Omega_Y^{m}$, và mặt khác ảnh của phần tử này trong
$H^{n-m}(X,\Omega_X^{n-m})$ chính là $P_X(Y)$. Những tính tương thích này,
vốn có thể được nêu ngay từ n\textsuperscript{o} 4, có thể được phát biểu, và thực ra đúng,
đối với mọi đa tạp không kỳ dị; tính tương thích thứ hai

\src{189}
được suy ra, chẳng hạn, từ tính giao hoán của sơ đồ các tự đồng cấu chính tắc
sau đây:
\begin{equation}\tag{7.6}
\resizebox{\textwidth}{!}{$
\left\{
\begin{array}{ccccc}
\operatorname{Ext}_{\Osh_X}^{n-m}(X;\Omega_X^{m},\Omega_X^{n})
&\longleftarrow&
\operatorname{Ext}_{\Osh_X}^{n-m}(X;\Omega_Y^{m},\Omega_X^{n})
&\simeq&
\operatorname{Hom}_{\Osh_X}(X;\Omega_Y^{m},\Omega_Y^{m})\\[3pt]
\Big\downarrow&&&&\Big\downarrow\\[3pt]
\operatorname{Ext}_{\Osh_X}^{n-m}(X;\Osh_X,\Omega_X^{n-m})
&\longleftarrow&
\operatorname{Ext}_{\Osh_X}^{n-m}(X;\Osh_Y,\Omega_X^{n-m})
&\simeq&
\operatorname{Hom}_{\Osh_X}(X;\Omega_X^{m},\Omega_Y^{m})
\end{array}
\right.
$}
\end{equation}

Như vậy ta thu được một bản tương tự chính xác của hình thức luận đối ngẫu
Poincaré trên các đa tạp compact có định hướng. Đặc biệt, Định lý 4 cho phép xác
định lớp đối đồng điều gắn với đường chéo của $X\times X$. Bằng một lập luận
quen thuộc, chẳng hạn ta suy ra một công thức Lefschetz:

\result{ĐỊNH LÝ}{5} Cho $f$ là một tự đồng cấu của một đa tạp xạ ảnh không
kỳ dị $X$, sao cho các điểm bất động của $f$ đều có bội 1. Khi đó, theo môđun
đặc số của $k$, số các điểm bất động này đồng dư với tổng luân phiên các vết
của các tự đồng cấu do $f$ cảm sinh trên các $H^{i}(X)$.

Hạn chế buộc phải đặt lên $f$ bắt nguồn từ những khó khăn được nêu trong
nhận xét ở n\textsuperscript{o} 4. Tuy nhiên, ta lưu ý rằng công thức Lefschetz
vẫn đúng nếu $f$ « đồng luân » với một tự đồng cấu mà mọi điểm bất động đều
có bội 1.

\fsection{8}{Khái quát hóa định lý đối ngẫu}

Cho $X$ là một đa tạp đại số không kỳ dị sao cho mọi bó đại số nhất quán
$F$ trên $X$ đều đẳng cấu với thương của một bó đại số nhất quán tự do địa
phương (điều này đúng nếu $X$ đóng địa phương trong một không gian xạ ảnh).
Khi đó mọi bó đại số nhất quán $F$ trên $X$ đều có một phân giải hữu hạn
$L$ bởi các bó tự do địa phương; với hai phân giải như vậy, ta luôn có thể
tìm được một phân giải thứ ba và một đồng cấu từ phân giải này vào từng phân
giải trong hai phân giải đầu, tương thích với các đồng cấu tăng cường. Tương
tự, nếu $L$ là một phân giải hữu hạn tự do địa phương của $F$, và nếu cho
một đồng cấu $F'\longrightarrow F$, thì tồn tại một phân giải hữu hạn tự do
địa phương $L'$ của $F'$ cùng một đồng cấu $L'\longrightarrow L$ tương thích
với $F'\longrightarrow F$; có thể giả sử đồng cấu này là toàn ánh nếu
$F'\longrightarrow F$ là toàn ánh. Điều này cho phép, với hai số nguyên
$r,s\geq0$ đã cho, định nghĩa hai đa hàm tử đối đồng điều theo các đối số
$A_1,\ldots,A_r$; $B_1,\ldots,B_s$,

\src{190}
biến thiên trong phạm trù các bó đại số nhất quán trên $X$, và các giá trị của
chúng lần lượt thuộc phạm trù các bó đại số nhất quán trên $X$ và phạm trù các
môđun trên $H^{0}(X,\Osh_X)$, theo các công thức:
\begin{equation}\tag{8.1}
\resizebox{\textwidth}{!}{$
\left\{
\begin{aligned}
\underline{T}_{r}^{s*}(A_1,\ldots,A_r;B_1,\ldots,B_s)
  &={}H^{*}\!\left(
    \Homsh_{\Osh_X}\!\left(
      L(A_1)\otimes\cdots\otimes L(A_r),
      L(B_1)\otimes\cdots\otimes L(B_s)
    \right)\right),\\[3pt]
T_{r}^{s*}(A_1,\ldots,A_r;B_1,\ldots,B_s)
  &={}\mathbf R^{*}\Gamma\!\left(
    \Homsh_{\Osh_X}\!\left(
      L(A_1)\otimes\cdots\otimes L(A_r),
      L(B_1)\otimes\cdots\otimes L(B_s)
    \right)\right).
\end{aligned}
\right.
$}
\end{equation}
Trong công thức này, $L(F)$ ký hiệu một phân giải hữu hạn tự do địa phương của
bó đại số nhất quán $F$, còn $\mathbf R^{*}\Gamma(K)$ ký hiệu siêu đối đồng
điều của không gian $X$ đối với phức các bó $K$. Nếu $r$ hoặc $s$ bằng 0, ta thay
tích tenxơ của các $L(A_i)$, tương ứng của các $L(B_j)$, bằng $\Osh_X$. Đặc
biệt, $\underline{T}_{0}^{0}$ và $T_{0}^{0}$ là các hàm tử phân bậc không có
đối số; $\underline{T}_{0}^{0}$ tập trung ở bậc 0, tại đó nó là bó $\Osh_X$,
còn $T_{0}^{0}$ được đồng nhất với $H^{*}(X,\Osh_X)$. Trong mọi trường hợp,
việc các vế phải của (8.1) không phụ thuộc vào các phân giải đã chọn là hiển
nhiên đối với dòng thứ nhất (vì khi đó vấn đề mang tính địa phương), còn đối
với dòng thứ hai thì suy ra từ những nhận xét tổng quát ở trên, khi xét dãy phổ
siêu đối đồng điều của phức các bó
\[
  K=\Homsh_{\Osh_X}\bigl(L(A_1)\otimes\cdots,
      L(B_1)\otimes\cdots\bigr),
\]
hội tụ tới siêu đối đồng điều của $X$ đối với $K$ và có số hạng ban đầu là
$H^{p}(X,H^{q}(K))$, tức là
\begin{equation}\tag{8.2}
  E_{2}^{p,q}
  =H^{p}\!\left(X,
    (\underline{T}_{r}^{s})^{(q)}
      (A_1,\ldots,A_r;B_1,\ldots,B_s)\right).
\end{equation}
Khi đó, ta thấy rằng chính dãy phổ này cũng không phụ thuộc vào các phân giải
đã chọn; đích hội tụ của nó là vế trái trong dòng thứ hai của (8.1). Ta dễ dàng
định nghĩa các toán tử nối ứng với từng đối số $A_i$, $B_j$, vì mọi dãy khớp
$0\longrightarrow F'\longrightarrow F\longrightarrow F''\longrightarrow0$
đều có thể được phân giải bởi một dãy khớp các phức hữu hạn tự do địa phương.

Trong hệ các hàm tử $\underline{T}_{r}^{s*}$, tương ứng $T_{r}^{s*}$, ta định
nghĩa các phép toán tương tự như các phép toán trong phép tính tenxơ; định nghĩa
của chúng suy ra ngay từ các công thức định nghĩa (8.1). Vì vậy, ta có một phép
hợp thành (khái quát hóa phép hợp thành đã xét ở n\textsuperscript{o} 2):
\begin{equation}\tag{8.3}
\resizebox{\textwidth}{!}{$
\begin{aligned}
&T_{r}^{s*}(A_1,\ldots,A_r;B_1,\ldots,B_s)
\times T_{r'}^{s'*}(A'_1,\ldots,A'_{r'};B'_1,\ldots,B'_{s'})\\
&\qquad\longrightarrow
T_{r+r'}^{s+s'*}
(A_1,\ldots,A_r,A'_1,\ldots,A'_{r'};
 B_1,\ldots,B_s,B'_1,\ldots,B'_{s'}),
\end{aligned}
$}
\end{equation}
thỏa mãn các tính chất hiển nhiên về tính kết hợp và tính tương thích với các
đồng cấu có tính hàm tử, các đồng cấu nối và các dãy phổ. Tương tự, ta có các
phép đối xứng, mà ta để độc giả tự viết tường minh. Ngoài ra, ta còn có một
phép co mỗi khi một trong các đối số $A_i$

\src{191}
trùng với một trong các đối số $B_j$:
\begin{equation}\tag{8.4}
\resizebox{\textwidth}{!}{$
\begin{aligned}
&T_{r}^{s*}(A_1,\ldots,A_{i-1},C,A_{i+1},\ldots,A_r;
 B_1,\ldots,B_{j-1},C,B_{j+1},\ldots,B_s)\\
&\qquad\longrightarrow
T_{r-1}^{s-1*}(A_1,\ldots,\widehat{A_i},\ldots,A_r;
 B_1,\ldots,\widehat{B_j},\ldots,B_s).
\end{aligned}
$}
\end{equation}
Ngoài ra, nếu một đối số $A_i$ là một bó tự do địa phương, thì có thể
lược bỏ đối số ấy, với điều kiện thay một trong các $A_j$ ($j\neq i$) bằng
$A_j\otimes A_i$, hoặc thay một trong các $B_k$ bằng $B_k\otimes A_i'$ (trong đó
$A_i'=\Homsh_{\Osh_X}(A_i,\Osh_X)$); ta có một quy tắc tính toán tương tự khi
một trong các đối số $B_i$ là một bó tự do địa phương. Đặc biệt, ta luôn có thể
lược bỏ một đối số đồng nhất với $\Osh_X$. Nếu mọi đối số đều tự do địa phương, trừ nhiều nhất
một trong các đối số $B_i$, thì quy tắc vừa nêu cho một đẳng cấu có tính hàm tử
\begin{equation}\tag{8.5}
  T_{r}^{s*}(A_1,\ldots,A_r;B_1,\ldots,B_s)
  =H^{*}\bigl(X,A_1'\otimes\cdots\otimes A_r'
    \otimes B_1\otimes\cdots\otimes B_s\bigr)
\end{equation}
(vì ta được quy về trường hợp $r=0$, $s=1$, trong đó điều này là hiển nhiên;
ta cũng có thể dùng trực tiếp dãy phổ có số hạng ban đầu (8.2)). Đối với các
$\underline{T}_{r}^{s}$, ta cũng định nghĩa các phép toán tương ứng với những
phép toán vừa nêu. Các quan hệ giữa những phép toán vừa được đưa vào cũng giống
như các quan hệ giữa những phép toán tương tự trong phép tính tenxơ.

Cho $n$ là chiều của $X$. Bằng cách áp dụng liên tiếp một phép hợp thành tenxơ
(8.3) và các phép co (8.4) trên những đối số trùng nhau, ta thu được một phép ghép cặp
\begin{equation}\tag{8.6}
\resizebox{\textwidth}{!}{$
  (T_{r}^{s})^{p}(A_1,\ldots;B_1,\ldots)
  \times
  (T_{s}^{r})^{n-p}(B_1,\ldots;A_1,\ldots,A_r\otimes\Omega_X^{n})
  \longrightarrow H^{n}(X,\Omega_X^{n}).
$}
\end{equation}

\result{ĐỊNH LÝ}{6} Nếu $X$ là một đa tạp xạ ảnh không kỳ dị, thì các phép
ghép cặp (8.6) là những phép đối ngẫu.

Điều này suy ra một cách hoàn toàn hình thức từ hệ quả của Định lý 3. Thật vậy,
từ hệ quả ấy dễ dàng suy ra rằng nếu $K$ là một phức các bó đại số nhất quán tự do
địa phương, thì siêu đối đồng điều của $X$ đối với $K$ đối ngẫu với siêu đối
đồng điều của $X$ đối với $K'\otimes\Omega_X^{n}$ qua các phép ghép cặp tự nhiên
\begin{equation}\tag{8.7}
  \mathbf R^{p}\Gamma(K)
  \times\mathbf R^{n-p}\Gamma(K'\otimes\Omega_X^{n})
  \longrightarrow\mathbf R^{n}\Gamma(\Omega_X^{n})
  =H^{n}(X,\Omega_X^{n}).
\end{equation}
Ta thấy điều này bằng cách dùng dãy phổ có số hạng ban đầu
$H^{p}(H^{q}(X,K))$ và dãy phổ tương tự cho
$K'\otimes\Omega_X^{n}$. Từ kết quả trên, Định lý 6 suy ra bằng cách dùng
định nghĩa (8.1).

\src{192}
\subsection*{NHẬN XÉT}

\noindent 1\textsuperscript{o} Đối với các định nghĩa trước Định lý 6,
không cần giả sử $X$ không kỳ dị, vì không nhất thiết chỉ làm việc với các
phân giải hữu hạn. Nhưng nếu $X$ kỳ dị, thì nói chung không còn có thể khẳng
định tiên nghiệm rằng các
$(\underline{T}_{r}^{s})^{p}(A_1,\ldots;B_1,\ldots)$ là những bó nhất quán,
bởi vì trong phức các bó
\[
  \Homsh_{\Osh_X}\bigl(L(A_1)\otimes\cdots,L(B_1)\otimes\cdots\bigr)
\]
sẽ có vô hạn thành phần có một bậc toàn phần cho trước.

\medskip
\noindent 2\textsuperscript{o} Dễ dàng kiểm tra rằng trong các công thức
(8.1), có thể thay một trong các $L(B_i)$ bằng $B_i$. Theo Mệnh đề 3, điều này
chứng tỏ rằng ta có
\begin{equation}\tag{8.8}
\left\{
\begin{aligned}
  \underline{T}_{1}^{1*}(A;B)&=\Extsh_{\Osh_X}^{*}(A,B),\\
  T_{1}^{1*}(A;B)&=\operatorname{Ext}_{\Osh_X}^{*}(X;A,B).
\end{aligned}
\right.
\end{equation}
Đặc biệt, lấy $r=s=1$ và $A_1=\Osh_X$ trong (8.6), ta thu lại Định lý 3.
Công thức (8.8) cũng kéo theo
\[
  T_{0}^{1*}(B)=H^{*}(X,B),
  \qquad
  T_{1}^{0*}(A)=\operatorname{Ext}_{\Osh_X}^{*}(X;A,\Osh_X).
\]

\medskip
\noindent 3\textsuperscript{o} Công thức (8.1) cho thấy rằng, nói chung, các
hàm tử $\underline{T}_{r}^{s*}$ và $T_{r}^{s*}$ có các thành phần ở cả bậc
dương lẫn bậc âm. Dùng Nhận xét 2, ta thấy rằng nếu $X$ có chiều $n$, thì các
thành phần khác không của $\underline{T}_{r}^{s*}$ nằm giữa $-(s-1)n$ và $rn$
khi $s>0$, và nằm giữa $0$ và $rn$ khi $s=0$; các thành phần khác không của
$T_{r}^{s*}$ nằm giữa $-(s-1)n$ và $(r+1)n$ khi $s>0$, và nằm giữa $0$ và
$(r+1)n$ khi $s=0$ (và thậm chí, nếu không nhầm, khi $r>0$, chúng nằm giữa
$-(s-1)n$ và $rn$, tương ứng giữa $0$ và $rn$).

\section*{TÀI LIỆU THAM KHẢO}

\noindent[1] Cartier (Pierre). --- Les groupes
$\operatorname{Ext}^{s}(A,B)$, Séminaire A. Grothendieck : Algèbre homologique,
tập 1, 1957, số 3.

\noindent[2] Grothendieck (Alexandre). --- Sur quelques points d'algèbre
homologique, \emph{Tohoku Math. J.}, tập 9, 1957, tr. 119--183.

\noindent[3] Serre (Jean-Pierre). --- Faisceaux algébriques cohérents,
\emph{Annals of Math.}, tập 61, 1955, tr. 197--278.

\noindent[4] Serre (Jean-Pierre). --- Sur la dimension homologique des anneaux
et des modules noethériens, \emph{Proc. Intern. Symp. on Alg. Number Theory}
[1955, Tokyo và Nikko]. --- Tokyo, Science Council of Japan, 1956;
tr. 175--189.

\clearpage

\src{193}
\section*{BỔ SUNG}

Những khó khăn được nêu trong nhận xét ở trang 13 hiện nay đã được giải quyết
hoàn toàn, nhờ việc mở rộng các định lý đối ngẫu sang các đa tạp bất kỳ
(có các kỳ dị tùy ý). Đối với lý thuyết này, vốn có thể được trình bày trong
khuôn khổ tổng quát của « lý thuyết lược đồ », chúng tôi dẫn chiếu đến bộ
« Éléments de Géométrie algébrique » của J. Dieudonné và A. Grothendieck,
hiện đang được biên soạn. Có thể tìm thấy một vài chỉ dẫn trong:

\begin{quote}
Grothendieck (A.). --- Cohomology theory of algebraic varieties,
\emph{Congrès int. des Mathématiciens}, 1958, Edinburgh (sắp xuất bản).
\end{quote}

\begin{flushright}
[Tháng 4 năm 1959]
\end{flushright}

\begin{center}\rule{36mm}{0.5pt}\end{center}
